\documentclass[11pt,a4paper]{scrartcl}
\usepackage{iutparakou}
\usepackage{tocloft}
\setlength{\parskip}{2pt}

\begin{document}

% ===================== PAGE DE GARDE =====================
\begin{titlepage}
\thispagestyle{empty}
\begin{tcolorbox}[enhanced,colback=iutblue,colframe=iutblue,sharp corners,
   coltext=white,height=6cm,valign=center,left=14pt,right=14pt]
  {\small IUT / IFRI — Université de Parakou}\\[2pt]
  {\small Master Génie Logiciel — Théorie des langages \& calculabilité}\\[14pt]
  {\Huge\bfseries De l'automate fini\\[4pt] à la machine universelle}\\[10pt]
  {\large\itshape Projet intégrateur — bibliothèque \texttt{formlang}}
\end{tcolorbox}

\vspace{10mm}
\begin{tcolorbox}[enhanced,colback=iutorange!8,colframe=iutorange,left=10pt,right=10pt]
\textbf{Énoncé étudiant} — cinq journées, une seule bibliothèque, quatre applications.\\[4pt]
\textbullet\ \textbf{Durée} : 1 semaine (5 jours + soutenance) \quad\textbullet\ \textbf{Binôme}
\quad\textbullet\ \textbf{Langage} : Python 3.10+ (\texttt{pytest} seul).\\[2pt]
\textbullet\ \textbf{Rendu} : dépôt Git \texttt{NOM1\_NOM2\_formlang/} + rapport PDF + soutenance.
\end{tcolorbox}

\vfill
\begin{center}
{\small Auteur du sujet : \textbf{Dr. Mêton ATINDEHOU}}\\[2pt]
{\footnotesize\color{iutgrey} Charte IUT/IFRI — bleu \texttt{\#1A3A6B} / orange \texttt{\#E07B20}}
\end{center}
\end{titlepage}

% ===================== SOMMAIRE =====================
\renewcommand{\cftsecfont}{\color{iutblue}\bfseries}
\tableofcontents
\clearpage

% ===================== CHAPEAU =====================
\section{Présentation générale}

\subsection{Le fil rouge}
Ce projet pose \textbf{une seule question}, à cinq niveaux de difficulté croissante :
\emph{reconnaître une structure} — depuis un simple motif dans un mot jusqu'à
n'importe quel calcul. À chaque étage, une leçon se répète :
\begin{quote}\itshape
plus la structure à reconnaître est riche, plus il faut de \textbf{mémoire} — et il
existe un seuil (la machine universelle) où l'on gagne l'universalité mais où l'on
\textbf{perd la décidabilité}.
\end{quote}
On ne construit pas « trois TP séparés » mais \textbf{un cœur unique} (\texttt{formlang/})
que l'on \textbf{réutilise} : deux applications instancient le même automate d'arbres, une
troisième partage la structure des mêmes termes, une quatrième décrit des tables de
machines de Turing exécutées par une machine universelle. Cette \emph{réutilisation}
est le cœur de l'évaluation.

\subsection{La hiérarchie de Chomsky, étage par étage}
\begin{center}\small
\begin{tabular}{@{}clll@{}}
\toprule
Jour & Étage & Pouvoir de reconnaissance & Mémoire \\\midrule
1 & Régulier & langages réguliers & état fini \\
2 & Hors-contexte & langages algébriques & pile (LIFO) \\
3 & \textbf{Arbres} & langages d'arbres réguliers & un état par nœud \\
4 & Calculable & récursivement énumérable & ruban libre \\
5 & Transversal & — & (congruence, intégration) \\
\bottomrule
\end{tabular}
\end{center}
\noindent Un mot $a_1\cdots a_n$ est un arbre dégénéré (branche unaire) ; un automate fini
est un automate d'arbres « tout en arité~1 ». Vous le \textbf{vérifierez en code}.

\subsection{Une bibliothèque, quatre applications}
Les applications n'ont pas le droit de réécrire un automate : elles \textbf{instancient}
ceux de \texttt{formlang}. C'est la règle « gate » la plus importante du projet.

\begin{center}
\begin{tikzpicture}[mod/.style={draw=iutblue,thick,rounded corners=2pt,fill=iutblue!8,
    font=\footnotesize\ttfamily,inner sep=3pt},
    app/.style={draw=iutorange!85!black,thick,rounded corners=2pt,fill=iutorange!12,
    font=\footnotesize\ttfamily,inner sep=3pt},
    fl/.style={-{Stealth[length=1.8mm]},thick,iutblue!70}]
  \node[mod] (tree) at (0,0) {formlang.tree};
  \node[mod] (dfa) at (0,-1.1) {formlang.dfa/nfa/fst/pda/cfg};
  \node[mod] (tur) at (0,-2.2) {formlang.turing $\to$ utm};
  \node[app] (morpho) at (-4,0.6) {apps/morpho};
  \node[app] (shield) at (-4,-0.4) {apps/shield};
  \node[app] (hash) at (-4,-1.4) {apps/hashcons};
  \node[app] (mtu) at (-4,-2.4) {apps/mtu};
  \node[app] (pipe) at (4,-1.1) {pipeline.py};
  \draw[fl] (morpho)--(tree); \draw[fl] (shield)--(tree); \draw[fl] (hash)--(tree);
  \draw[fl] (shield)--(dfa); \draw[fl] (mtu)--(tur);
  \draw[fl] (pipe)-- (morpho); \draw[fl] (pipe)-- (shield);
\end{tikzpicture}
\end{center}
\begin{description}[leftmargin=1.4em,style=nextline,font=\normalfont\ttfamily]
  \item[apps/morpho] morphologie — automate d'arbres (BUTA).
  \item[apps/shield] AttackDecomposer — automate d'arbres, 100\,\% structurel.
  \item[apps/hashcons] partage de structure (DAG) sur \texttt{formlang.tree.Term}.
  \item[apps/mtu] machine universelle $U$ + calculatrice unaire \emph{exécutée par} $U$.
\end{description}

\begin{ethique}
Le volet Shield porte \textbf{exclusivement} sur la \textbf{structure formelle} d'un
système de \emph{défense}. Tous les « prompts » sont des suites de \textbf{jetons
abstraits} ($a,o,r,e,[,],(,)$). \textbf{Aucune attaque réelle, aucun contenu sensible.}
Tout rendu contenant un exemple d'attaque réelle est hors-sujet et pénalisé.
\end{ethique}

\begin{gate}
\textbullet\ \emph{Anti-copier-coller} : une app qui \textbf{redéfinit} un automate au
lieu de l'instancier perd les points d'intégration de la journée.\\
\textbullet\ \emph{Honnêteté intellectuelle} : toute borne de complexité doit être
\textbf{démontrée ou référencée} (jamais « de mémoire »). Reportez vos \textbf{vrais}
chiffres de sortie.
\end{gate}

\subsection{Pré-requis}
Langages formels ; expressions régulières ; grammaires ; automates finis (AFD/AFN) ;
déterminisation ; minimisation ; notions de calculabilité. Le projet \textbf{introduit}
les automates d'arbres (comme généralisation des automates finis) et la machine
universelle (comme aboutissement).

\subsection{Organisation de la semaine \& barème global}
\begin{center}\small
\begin{tabular}{@{}cllr@{}}
\toprule
Jour & Thème & Tests cibles & Pts \\\midrule
1 & Régulier \& transducteurs & \texttt{test\_regular} & 4 \\
2 & Hors-contexte (pompage, CFG, PDA) & \texttt{test\_cfg\_nerode}, \texttt{test\_regular} & 4 \\
3 & \textbf{Arbres} + hash-consing & \texttt{test\_tree}, \texttt{test\_hashcons} & 5 \\
4 & Calculabilité (MT, MTU, calculatrice) & \texttt{test\_turing}, \texttt{test\_calc} & 4 \\
5 & Intégration \& Myhill--Nerode & \texttt{test\_pipeline}, \texttt{test\_cfg\_nerode} & 3 \\\midrule
\multicolumn{3}{@{}l}{\textbf{Total}} & \textbf{20} \\
\multicolumn{3}{@{}l}{Bonus (Thompson, CYK, minimisation BUTA, infixation, binaire)} & +6 \\
\bottomrule
\end{tabular}
\end{center}

\subsection{Méthode de travail (à lire avant de commencer)}
\textbullet\ Lancez \texttt{pytest -q} \textbf{immédiatement} : vous verrez $N$ échecs
(\texttt{NotImplementedError}). C'est normal.\\
\textbullet\ Chaque \texttt{\# TODO (Ex.)} du code renvoie à une question \textbf{E\dots}
d'une fiche-jour. Un exercice est \textbf{validé} quand son test passe au vert.\\
\textbullet\ Travaillez \emph{étage par étage} : ne passez au jour suivant que lorsque les
tests du jour courant sont verts.\\
\textbullet\ Le rapport se rédige \emph{en même temps} que le code (traces, preuves,
mesures), pas à la fin.

\subsection{Livrables}
\textbf{1.} Dépôt Git \texttt{NOM1\_NOM2\_formlang/} : \texttt{formlang/}, \texttt{apps/},
\texttt{pipeline.py}, \texttt{tests/} complétés et \textbf{qui passent} ; joindre
\texttt{pytest -q > out\_tests.txt} et \texttt{python pipeline.py > out\_demo.txt}.\\
\textbf{2.} Rapport PDF (gabarit \texttt{RAPPORT\_TEMPLATE.md}).\\
\textbf{3.} Soutenance (10 min) : démonstration bout-en-bout + un cas par étage.

\clearpage
% ===================== JOUR 1 =====================
\section{Jour 1 — Étage régulier : automates finis \& transducteurs}
\emph{Modules : \texttt{formlang/dfa.py}, \texttt{nfa.py}, \texttt{fst.py} — Apps :
\texttt{apps/shield/detector.py}, \texttt{normalizer.py} — Tests : \texttt{test\_regular.py}
— Barème : 4 pts (+1).}

\begin{ethique}Alphabet abstrait $\{a,o,r\}$ (+ chiffres \emph{leet}). Aucun contenu réel.\end{ethique}

\subsection{Contexte}
Le premier réflexe d'un détecteur est de repérer un \textbf{motif} : « le mot contient le
facteur \texttt{or} », un langage \textbf{régulier}. Avant de détecter, on
\textbf{normalise} (défaire le \emph{leetspeak} $4\!\to\!a$, $0\!\to\!o$, \dots) : une
relation \textbf{rationnelle} réalisée par un \textbf{transducteur fini}. On établit dès
ce jour la première \emph{limite due à la mémoire} (le renversement), qui annonce tout le
projet.

\subsection{Notions nouvelles}
\begin{objetformel}{l'automate fini}
\champ{Intuition.} Une machine à états qui lit un mot lettre par lettre, sans autre
mémoire que « où j'en suis » ; à la fin, elle répond \emph{oui} ou \emph{non}.

\champ{Définition formelle.} Un AFD est $M=(Q,\Sigma,\delta,q_0,F)$ avec
$\delta:Q\times\Sigma\to Q$ ; $M$ accepte $w$ si $\hat\delta(q_0,w)\in F$. Un langage est
\textbf{régulier} s'il est reconnu par un tel automate (de façon équivalente : par un AFN,
ou dénoté par une expression régulière).

\champ{Classe de problèmes.} Les \textbf{problèmes de reconnaissance de mots} d'un langage
régulier : décider l'appartenance $w\in L$. En pratique : recherche de motif, analyse
lexicale, validation de formats simples, filtres.

\champ{Dans le projet.} Décider « le mot contient le facteur \texttt{or} »
(\texttt{SingularityDetector}) — un pur problème d'appartenance à un langage régulier.
\end{objetformel}

\begin{objetformel}{le transducteur fini}
\champ{Intuition.} Un automate qui, en plus de lire, \emph{écrit} : il \textbf{traduit} un
mot en un autre au fil de la lecture.

\champ{Définition formelle.} Un transducteur séquentiel émet une sortie à chaque transition ;
il réalise une \textbf{relation rationnelle} (ici une fonction).

\champ{Classe de problèmes.} Les \textbf{transductions régulières} : normalisation,
translittération, remplacement, nettoyage de surface.

\champ{Dans le projet.} Normaliser le \emph{leetspeak} ($4\!\to\!a$, $0\!\to\!o$, \dots)
avant la détection.
\end{objetformel}

\subsubsection{AFD et AFN}
Un \textbf{AFD} lit un mot de gauche à droite ; dans chaque état, chaque lettre mène à
\emph{un seul} état. Un \textbf{AFN} autorise \emph{plusieurs} transitions pour une même
lettre : on accepte s'il \emph{existe} un chemin acceptant. L'AFN « contient \texttt{or} »
est très naturel : on « devine » où commence le facteur.
\begin{center}
\begin{tikzpicture}[node distance=24mm, on grid]
  \node[etat, initial, initial text={}] (q0) {$q_0$};
  \node[etat, right=of q0] (q1) {$q_1$};
  \node[etatfin, right=of q1, accepting] (q2) {$q_2$};
  \path[fleche]
    (q0) edge[loop above] node{$a,o,r$} (q0)
    (q0) edge node[above]{$o$} (q1)
    (q1) edge node[above]{$r$} (q2)
    (q2) edge[loop above] node{$a,o,r$} (q2);
\end{tikzpicture}
\end{center}
\noindent Non-déterminisme visible en $q_0$ : sur \texttt{o}, on peut \emph{rester} en
$q_0$ \textbf{ou} \emph{aller} en $q_1$.

\subsubsection{Déterminisation par construction des sous-ensembles}
Un état de l'AFD = un \emph{ensemble} d'états de l'AFN. On part de $\{q_0\}$ et on ferme
par transition :
\begin{center}\small
\begin{tabular}{@{}lccl@{}}
\toprule
état AFD & sur \texttt{o} & sur \texttt{a}, \texttt{r} & remarque\\\midrule
$S_0=\{q_0\}$ & $\{q_0,q_1\}$ & $\{q_0\}$ & initial\\
$S_1=\{q_0,q_1\}$ & $\{q_0,q_1\}$ & $\{q_0\}$/\,$\{q_0,q_2\}$ & \texttt o armé\\
$S_2=\{q_0,q_2\}$ & $\{q_0,q_1,q_2\}$ & $\{q_0,q_2\}$ & \textbf{accepteur}\\
\bottomrule
\end{tabular}
\end{center}
Après renommage et minimisation, on retrouve \textbf{trois} états — l'AFD minimal :
\begin{center}
\begin{tikzpicture}[node distance=26mm, on grid]
  \node[etat, initial, initial text={}] (A) {A};
  \node[etat, right=of A] (B) {B};
  \node[etatfin, right=of B, accepting] (C) {C};
  \path[fleche]
    (A) edge[loop above] node{$a,r$} (A)
    (A) edge[bend left] node[above]{$o$} (B)
    (B) edge[bend left] node[below]{$a$} (A)
    (B) edge[loop above] node{$o$} (B)
    (B) edge node[above]{$r$} (C)
    (C) edge[loop above] node{$a,o,r$} (C);
\end{tikzpicture}
\end{center}
\begin{exemple}[Trace de l'AFD sur \texttt{aor}]
\begin{center}\small
\begin{tabular}{@{}cccl@{}}
\toprule pas & état & lit & transition\\\midrule
0 & A & \texttt a & A $\xrightarrow{a}$ A\\
1 & A & \texttt o & A $\xrightarrow{o}$ B\\
2 & B & \texttt r & B $\xrightarrow{r}$ \textbf{C} (accepté)\\
\bottomrule
\end{tabular}
\end{center}
En \texttt B on est « armé » : un \texttt r immédiat bascule en \texttt C (absorbant) ;
\texttt{aoa} retombe de \texttt B en \texttt A, rejeté.
\end{exemple}

\subsubsection{Minimisation (Moore)}
Deux états sont \textbf{indistinguables} s'ils acceptent exactement les mêmes suffixes. On
part de la partition finaux/non-finaux et on \textbf{raffine} tant que deux états d'un bloc
mènent à des blocs différents. Le résultat (AFD minimal) est \textbf{unique} — lien direct
avec Myhill--Nerode (jour~5).

\subsubsection{Transducteur fini (FST) : normaliser le \emph{leet}}
Un AFD \emph{reconnaît} ; un \textbf{transducteur} \emph{traduit} (chaque transition émet
une sortie, notée $x\!:\!y$). Le normaliseur \emph{leet} tient en \textbf{un seul état}
(identité sur les lettres non listées).
\begin{center}
\begin{tikzpicture}[node distance=24mm, on grid]
  \node[etatfin, initial, initial text={}, accepting] (q) {$q_0$};
  \path[fleche] (q) edge[loop right] node[align=left]
     {$4\!:\!a\quad 0\!:\!o\quad 3\!:\!e$\\$1\!:\!i\quad 5\!:\!s\quad$(identité sinon)} (q);
\end{tikzpicture}
\end{center}
\begin{exemple}[Idempotence]
$T(\texttt{4or})=\texttt{aor}$, $T(\texttt{r0le})=\texttt{role}$. La sortie ne contient
plus de chiffre \emph{leet}, donc $T(T(w))=T(w)$ : \textbf{idempotent}. La composition
$T_2\circ T_1$ se lit sur les paires d'états (\texttt{compose}).
\end{exemple}

\subsubsection{Une limite déjà visible : one-way vs two-way}
Le renversement $w\mapsto w^{R}$ est \textbf{impossible} en \textbf{one-way} : il faudrait
mémoriser tout $w$ avant d'émettre la première lettre — impossible à mémoire bornée. Une
tête \textbf{two-way} le permet. C'est le premier exemple concret de \emph{limite due à la
mémoire} — thème qui culminera au jour~4.

\subsection{A. Questions théoriques \normalfont(à rédiger) — 2 pts}
On pose $\Sigma=\{a,o,r\}$ et $L_1=\{w \mid w \text{ contient le facteur } or\}$.
\begin{description}[leftmargin=1.8em,style=nextline]
  \item[Q1.1 (0,5) — expression régulière.] Donner une regex dénotant \emph{exactement}
    $L_1$. \emph{Livrable} : la regex + une phrase justifiant qu'elle capture tous les mots
    contenant \texttt{or} et aucun autre.
  \item[Q1.2 (1) — du non-déterminisme à l'AFD minimal.] Construire un AFN pour $L_1$, le
    \textbf{déterminiser} (construction des sous-ensembles) puis le \textbf{minimiser}
    (raffinement de Moore). \emph{Livrables} : (a) l'AFN dessiné ; (b) la \textbf{table
    complète} des sous-ensembles atteints ; (c) la \textbf{suite des partitions} de
    minimisation ; (d) l'AFD minimal dessiné + son nombre d'états.
    \emph{Indication} : vous devez retrouver \textbf{3 états}.
  \item[Q1.3 (0,5) — facteur vs sous-séquence.] Soit $L_1'=\{w\mid w$ contient un \texttt o
    suivi, \emph{pas forcément immédiatement}, d'un \texttt r$\}$. Donner sa regex,
    expliquer la différence avec $L_1$, et prouver l'inclusion $L_1\subseteq L_1'$.
    \emph{Livrable} : regex + argument d'inclusion + un mot de $L_1'\setminus L_1$.
\end{description}

\subsection{B. Exercices de programmation — 2 pts}
\begin{description}[leftmargin=1.8em,style=nextline]
  \item[E1.1 (0,75) — détecteur AFD.] Dans \texttt{apps/shield/detector.py}, compléter la
    table $\delta$ de \texttt{\_DFA\_OR} (états \texttt{A},\texttt{B},\texttt{C}) et, dans
    \texttt{formlang/dfa.py}, \texttt{DFA.run}/\texttt{DFA.accepts}. \emph{Contrainte} :
    \texttt{contains\_or} passe \textbf{par l'AFD}, jamais par \texttt{"or" in w}.
    \emph{Test} : \texttt{test\_detector\_or}.
  \item[E1.2 (0,5) — minimisation.] \texttt{DFA.minimize} (raffinement de Moore ; utiliser
    \texttt{\_completed()} fourni). \emph{Test} : \texttt{test\_minimisation\_3\_etats}.
  \item[E1.3 (0,5) — déterminisation.] \texttt{NFA.to\_dfa} (sous-ensembles ;
    \texttt{\_eps\_closure}/\texttt{\_move} fournis). \emph{Test} :
    \texttt{test\_nfa\_to\_dfa\_equivalence}.
  \item[E1.4 (0,25) — transducteur.] \texttt{leet\_fst} (un état final,
    \texttt{identity\_on\_missing=True}) + \texttt{SequentialFST.transduce}. \emph{Tests} :
    \texttt{test\_leet\_idempotent}, \texttt{test\_fst\_composition}.
\end{description}

\begin{notion}[Pièges fréquents]
\textbullet\ Oublier que \texttt{DFA.run} doit renvoyer \texttt{None} si une transition
manque (mot rejeté), pas lever d'exception.\\
\textbullet\ Minimisation : ne pas compléter l'automate (états puits) avant de raffiner
$\Rightarrow$ partitions fausses.\\
\textbullet\ FST : émettre le caractère \emph{lu} plutôt que la \emph{sortie} sur les
transitions listées.
\end{notion}

\begin{exemple}[Sortie de référence]\ttfamily
\$ pytest -q tests/test\_regular.py\\ 7 passed
\end{exemple}

\begin{center}\small
\begin{tabular}{@{}llr@{}}
\toprule Item & Détail & Pts\\\midrule
Q1.1–Q1.3 & regex, AFN$\to$AFD$\to$min, facteur vs sous-séquence & 2\\
E1.1 & détecteur AFD (\texttt{contains\_or}) & 0,75\\
E1.2 & minimisation 3 états & 0,5\\
E1.3 & déterminisation & 0,5\\
E1.4 & FST leet + composition & 0,25\\
\textbf{Bonus} & construction de Thompson (regex $\to$ AFN) & +1\\
\bottomrule
\end{tabular}
\end{center}
\begin{gate}\texttt{detector.py} \textbf{instancie} \texttt{formlang.dfa.DFA} ; pas de \texttt{"or" in w}.\end{gate}
\clearpage
% ===================== JOUR 2 =====================
\section{Jour 2 — Étage hors-contexte : grammaires \& automate à pile}
\emph{Modules : \texttt{formlang/pda.py}, \texttt{cfg.py} — App : \texttt{apps/shield/delimiters.py}
— Tests : \texttt{test\_cfg\_nerode.py}, \texttt{test\_regular.py} — Barème : 4 pts (+1).}
\begin{ethique}$\Sigma_3=\{a,o,r,[,],(,)\}$. Aucun contenu réel.\end{ethique}

\subsection{Contexte}
Les \textbf{faux délimiteurs} imitent un bloc \texttt{[\,\dots\,]} ou un cadre
\texttt{(\,\dots\,)}, éventuellement imbriqués. Vérifier l'imbrication bien formée =
reconnaître un langage de type \textbf{Dyck} : c'est \textbf{hors-contexte}, hors de portée
d'un automate fini. Le nouvel ingrédient : la \textbf{pile}.

\subsection{Notions nouvelles}
\begin{objetformel}{la grammaire hors-contexte \& l'automate à pile}
\champ{Intuition.} Une grammaire \emph{engendre} des structures imbriquées par réécriture ;
l'automate à pile est un automate fini \textbf{muni d'une pile} pour se souvenir de la
profondeur d'imbrication.

\champ{Définition formelle.} Une CFG est $G=(N,\Sigma,R,S)$ ; un PDA ajoute une pile à un
automate fini. Les deux décrivent la \textbf{même} classe : les langages \textbf{algébriques}
(hors-contexte).

\champ{Classe de problèmes.} La \textbf{reconnaissance et l'analyse syntaxique} de langages
hors-contexte : équilibrage de parenthèses, expressions arithmétiques, structures
récursives, \emph{parsing} de langages de programmation.

\champ{Dans le projet.} Vérifier l'imbrication bien formée des délimiteurs \texttt{[\,]},
\texttt{(\,)} — impossible pour un automate fini (jour~1), naturel pour une pile.
\end{objetformel}

\subsubsection{Le lemme de l'étoile (pompage)}
Outil pour \textbf{prouver qu'un langage n'est pas régulier}. Si $L$ est régulier, il
existe $p$ tel que tout mot long $w=xyz$ ($|xy|\le p$, $y\neq\varepsilon$) vérifie
$xy^{i}z\in L$ pour tout $i$. Pour $D=\{[^{n}]^{n}\}$, prendre $w=[^{p}]^{p}$ : comme
$|xy|\le p$, le bloc $y$ n'est fait que de \texttt{[}, donc $xy^{2}z=[^{p+k}]^{p}\notin D$.
\begin{center}
\begin{tikzpicture}[x=6.5mm]
  \foreach \i in {0,...,3} {\node[draw=iutblue,thick,minimum size=6mm,font=\ttfamily] at (\i,0) {[};}
  \foreach \i in {4,...,7} {\node[draw=iutblue,thick,minimum size=6mm,font=\ttfamily] at (\i,0) {]};}
  \draw[decorate,decoration={brace,amplitude=4pt},iutorange!85!black,thick]
     (-0.35,0.45) -- (1.35,0.45) node[midway,above=3pt,font=\scriptsize]{$x$};
  \draw[decorate,decoration={brace,amplitude=4pt},iutorange!85!black,thick]
     (1.65,0.45) -- (2.35,0.45) node[midway,above=3pt,font=\scriptsize]{$y=[^{k}$};
  \draw[decorate,decoration={brace,amplitude=4pt},iutorange!85!black,thick]
     (2.65,0.45) -- (7.35,0.45) node[midway,above=3pt,font=\scriptsize]{$z$};
\end{tikzpicture}
\end{center}
Un automate fini a une mémoire \textbf{bornée} : il ne peut pas « compter » les ouvrantes.

\subsubsection{Grammaire hors-contexte et arbre de dérivation}
Une CFG engendre par \textbf{réécriture} :
$S \to S\,S \mid [\,S\,] \mid (\,S\,) \mid a \mid o \mid r \mid \varepsilon.$
L'arbre de dérivation montre \emph{comment} un mot est produit (\texttt{[a]} et \texttt{([])}) :
\begin{center}
\begin{tikzpicture}[level distance=10mm, sibling distance=9mm,
  every node/.style={noeud, align=center, inner sep=2pt, font=\footnotesize},
  edge from parent/.style={draw=iutblue!70, thick}, baseline=(current bounding box.north)]
\node {S} child { node[feuille] {[} } child { node {S} child { node[feuille] {a} } } child { node[feuille] {]} };
\end{tikzpicture}\qquad\qquad
\begin{tikzpicture}[level distance=10mm, sibling distance=9mm,
  every node/.style={noeud, align=center, inner sep=2pt, font=\footnotesize},
  edge from parent/.style={draw=iutblue!70, thick}, baseline=(current bounding box.north)]
\node {S} child { node[feuille] {(} }
  child { node {S} child { node[feuille] {[} } child { node {S} child { node[feuille] {$\varepsilon$} } } child { node[feuille] {]} } }
  child { node[feuille] {)} };
\end{tikzpicture}
\end{center}

\subsubsection{Ambiguïté}
Une grammaire est \textbf{ambiguë} si un mot admet plusieurs arbres. $S\to SS$ n'est pas
associée : \texttt{aaa} a (au moins) deux arbres.
\begin{center}
\begin{tikzpicture}[level distance=9mm, sibling distance=8mm,
  every node/.style={noeud, align=center, inner sep=2pt, font=\footnotesize},
  edge from parent/.style={draw=iutblue!70, thick}, baseline=(current bounding box.north)]
\node {S} child { node {S} child { node[feuille]{a} } }
  child { node {S} child { node {S} child { node[feuille]{a} } } child { node {S} child { node[feuille]{a} } } };
\end{tikzpicture}\qquad\qquad
\begin{tikzpicture}[level distance=9mm, sibling distance=8mm,
  every node/.style={noeud, align=center, inner sep=2pt, font=\footnotesize},
  edge from parent/.style={draw=iutblue!70, thick}, baseline=(current bounding box.north)]
\node {S} child { node {S} child { node {S} child { node[feuille]{a} } } child { node {S} child { node[feuille]{a} } } }
  child { node {S} child { node[feuille]{a} } };
\end{tikzpicture}
\end{center}
On désambiguïse par $S\to T\,S \mid T$ ; $T\to [\,S\,]\mid(\,S\,)\mid a\mid o\mid r\mid\varepsilon$.

\subsubsection{La pile : mémoire non bornée mais disciplinée}
La \textbf{pile} (LIFO) est \textbf{non bornée} mais accessible seulement au \textbf{sommet}.
Évolution sur \texttt{[a(r)]} (on ignore \texttt a, \texttt r) : accepté car pile vide à la fin.
\begin{center}
\begin{tikzpicture}[pile/.style={draw=iutblue,thick,minimum width=8mm,minimum height=6mm,font=\ttfamily},
    lab/.style={font=\scriptsize\color{iutgrey}}, x=17mm]
  \node[pile,fill=iutblue!8] at (0,0) {[}; \node[lab] at (0,-0.55){lit \texttt{[}};
  \node[pile,fill=iutblue!8] at (1,0) {[}; \node[pile,fill=iutorange!15] at (1,0.6) {(};
  \node[lab] at (1,-0.55){lit \texttt{(}};
  \node[pile,fill=iutblue!8] at (2,0) {[}; \node[lab] at (2,-0.55){lit \texttt{)} : dépile};
  \node[draw=iutblue!40,dashed,minimum width=8mm,minimum height=6mm] at (3,0){};
  \node[lab] at (3,-0.55){lit \texttt{]} : vide};
  \node[iutorange!85!black,font=\small\bfseries] at (3.95,0){ACCEPTÉ};
  \draw[-{Stealth[length=2mm]},iutblue!70,thick] (0.35,0)--(0.65,0);
  \draw[-{Stealth[length=2mm]},iutblue!70,thick] (1.35,0)--(1.65,0);
  \draw[-{Stealth[length=2mm]},iutblue!70,thick] (2.35,0)--(2.65,0);
\end{tikzpicture}
\end{center}
\begin{exemple}[Un rejet : \texttt{([)]}]
On empile \texttt( puis \texttt[. Sur \texttt), le sommet est \texttt[ $\neq$ ouvrante
attendue \texttt( $\Rightarrow$ \textbf{rejet}. L'imbrication « croisée » est mal formée :
c'est la \emph{pile} qui l'attrape, pas un compteur.
\end{exemple}

\subsection{A. Questions théoriques \normalfont(à rédiger) — 2,5 pts}
\begin{description}[leftmargin=1.8em,style=nextline]
  \item[Q2.1 (1,5) — non-régularité par pompage.] Montrer rigoureusement que
    $D=\{[^{n}]^{n}\mid n\ge 0\}$ n'est pas régulier. \emph{Livrable} : la preuve complète
    (choix de $w$, discussion de \emph{toute} découpe $xyz$ admissible, mot pompé hors de
    $D$), puis la conclusion : le \texttt{SingularityDetector} (AFD) ne peut pas vérifier
    l'imbrication.
  \item[Q2.2 (0,5) — grammaire.] Donner une grammaire $G$ engendrant les prompts bien
    parenthésés sur $\Sigma_3$ (\texttt a, \texttt o, \texttt r libres entre délimiteurs).
    \emph{Livrable} : les règles + un exemple de dérivation.
  \item[Q2.3 (0,5) — ambiguïté.] $G$ est-elle ambiguë ? Si oui, donner \textbf{deux} arbres
    de dérivation de \texttt{aaa}, puis une grammaire \textbf{non ambiguë} équivalente.
\end{description}

\subsection{B. Exercices de programmation — 1,5 pt}
\begin{description}[leftmargin=1.8em,style=nextline]
  \item[E2.1 (1) — automate à pile.] \texttt{DelimiterPDA.accepts} \emph{avec une pile}
    (liste) : empiler sur ouvrant, dépiler sur fermant si le sommet correspond, ignorer
    \texttt{a,o,r,e}, accepter si pile vide. \emph{Test} : \texttt{test\_delimiters\_pda}.
  \item[E2.2 (0,5) — génération bornée.] \texttt{CFG.generate(max\_len)} : énumérer les mots
    \emph{terminaux} de longueur $\le$ \texttt{max\_len}. \emph{Test} :
    \texttt{test\_cfg\_engendre\_des\_mots\_equilibres}.
\end{description}
\begin{notion}[Piège majeur de E2.2]
Avec $S\to SS\mid\dots\mid\varepsilon$, des formes purement non-terminales (\texttt{S},
\texttt{SS}, \texttt{SSS}, \dots) ont \textbf{0 terminal} et ne sont jamais élaguées par la
seule borne sur les terminaux $\Rightarrow$ \textbf{explosion mémoire}. Il faut borner
\textbf{aussi} le nombre de non-terminaux (p.\,ex.\ $\le 2\cdot\texttt{max\_len}+2$) ; cette
borne ne retire aucun mot de longueur $\le$ \texttt{max\_len}.
\end{notion}
\begin{center}\small
\begin{tabular}{@{}llr@{}}
\toprule Item & Détail & Pts\\\midrule
Q2.1 & pompage $\{[^{n}]^{n}\}$ & 1,5\\
Q2.2 & grammaire bien parenthésée & 0,5\\
Q2.3 & ambiguïté + désambiguïsation & 0,5\\
E2.1 & PDA \texttt{well\_parenthesized} & 1\\
E2.2 & génération bornée & 0,5\\
\textbf{Bonus} & reconnaissance par \textbf{CYK} & +1\\
\bottomrule
\end{tabular}
\end{center}
\begin{gate}\texttt{delimiters.py} \textbf{instancie} \texttt{formlang.pda.DelimiterPDA}.\end{gate}
\clearpage
% ===================== JOUR 3 =====================
\section{Jour 3 — Automates d'arbres (pivot) \& hash-consing}
\emph{Module : \texttt{formlang/tree.py} — Apps : \texttt{apps/morpho}, \texttt{apps/shield},
\texttt{apps/hashcons} — Tests : \texttt{test\_tree.py}, \texttt{test\_hashcons.py} —
Barème : 5 pts (+2).}
\begin{ethique}Volet Shield 100\,\% structurel : jetons abstraits, aucune attaque réelle.\end{ethique}

\subsection{Contexte}
C'est le \textbf{pivot} du projet. Une séquence plate \texttt{p p r s} ne dit pas
\emph{comment} le mot est construit : la morphologie est \textbf{hiérarchique}, donc un
\textbf{arbre}. On généralise l'automate fini : au lieu de suivre une \emph{chaîne}
d'états, on \textbf{fusionne} les états des enfants d'un nœud. Deux applications très
différentes (morphologie, pare-feu structurel) \emph{et} le hash-consing instancient le
\textbf{même} \texttt{formlang.tree}.

\subsection{Notions nouvelles}
\begin{objetformel}{l'automate d'arbres}
\champ{Intuition.} Un automate fini qui, au lieu de parcourir une \emph{ligne}, « lit un
\textbf{arbre} » des feuilles vers la racine, et décide si sa \textbf{forme} est admissible.

\champ{Définition formelle.} Un BUTA est $A=(Q,\Sigma,Q_f,\Delta)$ avec des règles
$f(q_1,\dots,q_n)\to q$ ; il accepte si la racine reçoit un état de $Q_f$. Il reconnaît les
\textbf{langages d'arbres réguliers} (généralisation directe des langages réguliers de mots :
un mot est un arbre unaire).

\champ{Classe de problèmes.} La \textbf{validation de structures hiérarchiques} : conformité
de schémas XML/JSON, vérification d'arbres de syntaxe abstraite (AST), typage, filtrage
d'arbres (\emph{tree matching}), \emph{model checking} d'arbres.

\champ{Dans le projet.} Valider la structure \emph{morphologique} d'un mot (\texttt{morpho})
\emph{et} décomposer/bloquer des structures dangereuses (\texttt{shield}) — deux applications
d'un \textbf{même} automate d'arbres.
\end{objetformel}

\begin{objetformel}{le hash-consing (partage de structure)}
\champ{Intuition.} Ne ranger chaque sous-arbre distinct \textbf{qu'une seule fois} et
partager les copies : l'arbre devient un graphe orienté acyclique (\textbf{DAG}).

\champ{Définition formelle.} Une fonction d'\emph{internement} injective (terme $\mapsto$
identifiant) via une table canonique, garantissant un \textbf{partage maximal} et l'égalité
structurelle en $O(1)$.

\champ{Classe de problèmes.} La \textbf{représentation compacte} et la comparaison rapide de
structures : dictionnaires morphologiques minimaux, diagrammes de décision (BDD), partage de
sous-expressions en compilation.

\champ{Dans le projet.} Mesurer le partage et le \textbf{taux de compression} des arbres
morphologiques d'un vocabulaire.
\end{objetformel}

\subsubsection{Alphabet classé, termes, et automate d'arbres ascendant (BUTA)}
Chaque symbole a une \textbf{arité}. Un \textbf{BUTA} est $A=(Q,\Sigma,Q_f,\Delta)$ avec
des règles $f(q_1,\dots,q_n)\to q$. On accepte si la \textbf{racine} reçoit un état de $Q_f$.
\begin{notion}[L'AFD généralisé]
Un AFD lit un mot de gauche à droite : cas « tout en arité~1 » (chaîne). Le BUTA part des
\textbf{feuilles} et \textbf{remonte} : à un nœud, il \emph{combine} les états des enfants.
\textbullet\ Code : \texttt{TreeAutomaton.run} (post-ordre ; \texttt{REJECT} si un enfant
rejette ou si aucune règle).
\end{notion}

\subsubsection{Exécution ascendante — exemple morphologique}
Mot circonfixé \texttt{a-na-pend-a} : l'état de chaque nœud (\textcolor{iutorange!85!black}{orange})
est calculé de bas en haut ; racine \textcolor{iutorange!85!black}{\texttt{WORD}}$\in Q_f$
$\Rightarrow$ accepté (classe \textsc{circumfixed}).
\begin{center}
\begin{tikzpicture}[level distance=14mm,
  level 1/.style={sibling distance=50mm}, level 2/.style={sibling distance=25mm},
  level 3/.style={sibling distance=21mm}, every node/.style={noeud, align=center},
  edge from parent/.style={draw=iutblue!70, thick}]
\node {word\\{\scriptsize\bfseries\color{iutorange!85!black}WORD}}
  child { node {prefixes\\{\scriptsize\bfseries\color{iutorange!85!black}PREFS}}
    child { node[feuille] {prefix\,\texttt{a}\\{\scriptsize\bfseries\color{iutorange!85!black}PRE}} }
    child { node {prefixes\\{\scriptsize\bfseries\color{iutorange!85!black}PREFS}}
      child { node[feuille] {prefix\,\texttt{na}\\{\scriptsize\bfseries\color{iutorange!85!black}PRE}} }
      child { node[feuille] {nil\\{\scriptsize\bfseries\color{iutorange!85!black}NIL}} } } }
  child { node {rest\\{\scriptsize\bfseries\color{iutorange!85!black}REST}}
    child { node[feuille] {root\,\texttt{pend}\\{\scriptsize\bfseries\color{iutorange!85!black}ROOT}} }
    child { node {suffixes\\{\scriptsize\bfseries\color{iutorange!85!black}SUFS}}
      child { node[feuille] {suffix\,\texttt{a}\\{\scriptsize\bfseries\color{iutorange!85!black}SUF}} }
      child { node[feuille] {nil\\{\scriptsize\bfseries\color{iutorange!85!black}NIL}} } } };
\end{tikzpicture}
\end{center}

\subsubsection{Pourquoi un AFD \emph{de mots} échoue : la frontière ne suffit pas}
La \textbf{frontière} (yield) ne détermine pas la structure. Les deux arbres ci-dessous ont
la \textbf{même} frontière \texttt{na pend} mais des classes différentes ; un AFD, qui ne
lit que la frontière, est \textbf{aveugle} à la différence.
\begin{center}
\begin{tikzpicture}[level distance=11mm, sibling distance=15mm,
  every node/.style={noeud, align=center, font=\scriptsize, inner sep=2pt},
  edge from parent/.style={draw=iutblue!70, thick}, baseline=(current bounding box.north)]
\node {word} child { node {prefixes} child { node[feuille] {prefix \texttt{na}} } }
  child { node {rest} child { node[feuille] {root \texttt{pend}} } };
\end{tikzpicture}\hspace{3mm}{\small\itshape(PREFIXED)}\qquad
\begin{tikzpicture}[level distance=11mm, sibling distance=15mm,
  every node/.style={noeud, align=center, font=\scriptsize, inner sep=2pt},
  edge from parent/.style={draw=iutblue!70, thick}, baseline=(current bounding box.north)]
\node {word} child { node {rest} child { node[feuille] {root \texttt{na}} }
    child { node {suffixes} child { node[feuille] {suffix \texttt{pend}} } } };
\end{tikzpicture}\hspace{3mm}{\small\itshape(SUFFIXED)}
\end{center}
\begin{notion}[Théorème du yield — pont vers le jour 1]
La frontière d'un mot morphologique accepté est exactement $p^{*}\,r\,s^{*}$, le langage
\textbf{régulier} du jour~1 ; réciproquement, l'ensemble des frontières d'un automate
d'arbres régulier est \textbf{hors-contexte}. Déterminisme $\Rightarrow$ exécution unique,
coût $O(n)$.
\end{notion}

\subsubsection{Le même moteur pour le Shield (100\,\% structurel)}
\texttt{sys(seq(txt, frame(role)))} remonte à \textbf{DANGER}$\in Q_f$ $\Rightarrow$ bloqué.
Que de la \emph{forme}.
\begin{center}
\begin{tikzpicture}[level distance=13mm,
  level 1/.style={sibling distance=44mm}, level 2/.style={sibling distance=22mm},
  every node/.style={noeud, align=center, font=\footnotesize},
  edge from parent/.style={draw=iutblue!70, thick}]
\node {sys\\{\scriptsize\bfseries\color{red!70!black}DANGER}}
  child { node {seq\\{\scriptsize\bfseries\color{red!70!black}DANGER}}
    child { node[feuille] {txt\\{\scriptsize\bfseries\color{iutorange!85!black}SAFE}} }
    child { node {frame\\{\scriptsize\bfseries\color{red!70!black}DANGER}}
      child { node[feuille] {role\\{\scriptsize\bfseries\color{iutorange!85!black}ROLE}} } } };
\end{tikzpicture}
\end{center}

\subsubsection{Clôture par intersection (produit) et hash-consing (DAG)}
Clôture par \textbf{intersection} : $A_1\times A_2$ a pour états des \emph{paires}, règles
composante par composante. Le \textbf{hash-consing} donne à chaque sous-arbre distinct un
\textbf{id unique} : sous-arbres identiques \textbf{partagés}, l'arbre devient un \textbf{DAG}.
\begin{center}
\begin{tikzpicture}[every node/.style={noeud, font=\scriptsize, inner sep=2pt},
   fl/.style={-{Stealth[length=1.6mm]}, thick, draw=iutblue!70}]
  \node (w1) at (0,0) {word 1}; \node (w2) at (3.4,0) {word 2};
  \node[feuille] (root1) at (-0.9,-1.2) {root \texttt{struct}};
  \node[feuille] (pre) at (2.4,-1.2) {prefix \texttt{re}};
  \node[noeud, fill=iutorange!15, draw=iutorange!85!black] (shared) at (1.7,-2.4) {suffixes(\texttt{ation}, nil)};
  \draw[fl] (w1)--(root1); \draw[fl] (w2)--(pre); \draw[fl] (w1)--(shared); \draw[fl] (w2)--(shared);
  \node[font=\scriptsize\itshape, text=iutgrey, text width=44mm] at (7.0,-1.2)
    {Le sous-arbre \texttt{-ation} n'est stocké \textbf{qu'une fois}. round-trip :
     \texttt{get(intern(t))==t} ; compression $1-\text{uniques}/\text{total}$.};
\end{tikzpicture}
\end{center}

\subsection{A. Questions théoriques \normalfont(à rédiger) — 2 pts}
\begin{description}[leftmargin=1.8em,style=nextline]
  \item[Q3.1 (0,5) — déterminisme \& coût.] Montrer que \texttt{morpho\_automaton} et
    \texttt{shield\_automaton} sont déterministes (au plus une règle par
    $(\text{symbole},\text{états enfants})$). \emph{Livrable} : l'argument + la conséquence
    (unicité de l'exécution, coût $O(n)$).
  \item[Q3.2 (0,5) — la frontière ne suffit pas.] Donner deux arbres de \textbf{même
    frontière} et de classes différentes, et expliquer pourquoi un AFD de mots ne peut les
    distinguer.
  \item[Q3.3 (0,5) — théorème du yield.] Montrer que la frontière d'un mot morphologique
    accepté est $p^{*}r s^{*}$, et faire le lien explicite avec le langage régulier du jour~1.
  \item[Q3.4 (0,5) — réduplication.] La réduplication (copie d'une partie de la racine)
    engendre-t-elle un langage d'arbres \emph{régulier} ? Argumenter via $\{ww\}$ (non
    régulier). \emph{Indice réel} : Culy 1985 — le bambara sort du hors-contexte.
\end{description}

\subsection{B. Exercices de programmation — 3 pts}
\begin{description}[leftmargin=1.8em,style=nextline]
  \item[E3.1 (1) — moteur BUTA.] \texttt{TreeAutomaton.run} (post-ordre ; \texttt{REJECT} si
    enfant rejeté ou règle absente) + \texttt{accepts}. \emph{Débloque tout
    \texttt{test\_tree.py}.}
  \item[E3.2 (0,5) — automate morphologique.] Règles $\Delta$ de \texttt{morpho\_automaton}
    (feuilles \texttt{prefix/root/suffix/nil} ; chaînes ; \texttt{rest}, \texttt{word}).
    \emph{Tests} : \texttt{test\_morpho\_accept\_et\_classes}, \texttt{test\_morpho\_rejet}.
  \item[E3.3 (0,5) — AttackDecomposer.] \texttt{\_seq} (DANGER si l'un est DANGER ou si
    $\{$OVR,ROLE$\}$, sinon sévérité max) + \texttt{shield\_automaton}. \emph{Test} :
    \texttt{test\_shield\_blocage}.
  \item[E3.4 (0,5) — produit.] \texttt{product(A1,A2)} : $L=L(A_1)\cap L(A_2)$ (états
    = paires, règles composante par composante). \emph{Test} : \texttt{test\_produit\_intersection}.
  \item[E3.5 (0,5) — hash-consing.] \texttt{CompactStore.intern} (partage) + \texttt{get}
    (round-trip). \emph{Tests} : \texttt{test\_round\_trip\_exact},
    \texttt{test\_sous\_arbres\_identiques\_partages}, \texttt{test\_compression\_positive}.
\end{description}
\textbf{Q3.5 (rapport)} : $\text{compression}=1-\text{uniques}/\text{total}$ ; discuter pour
quelle famille de langues (anglais isolant vs turc/finnois agglutinants) le partage est le
plus élevé. Reportez vos \textbf{vrais} chiffres.

\textbf{E3.6 (bonus) — P4.5 : « dangereux \emph{et} doublement encodé ».} Construire \texttt{enc\_automaton} (compte les feuilles \texttt{enc}, plafonné à 2, final $=\{2\}$, même signature), puis \texttt{dangerous\_and\_double\_encoded = product(shield\_automaton(), enc\_automaton())}. On bloque \emph{ssi} l'arbre est dangereux \textbf{et} porte $\ge 2$ encodages : illustration directe de la \textbf{clôture par intersection}. \emph{Test} : \texttt{test\_shield\_double\_encodage\_P45}.

\begin{notion}[Pièges fréquents]
\textbullet\ \texttt{run} : oublier de propager \texttt{REJECT} dès qu'un enfant rejette.\\
\textbullet\ \texttt{shield} : oublier une paire $(x,y)$ dans les règles \texttt{seq}
(le produit cartésien des états).\\
\textbullet\ \texttt{intern} : clé canonique incomplète (oublier \texttt{label}) $\Rightarrow$
faux partage et round-trip cassé.
\end{notion}
\begin{center}\small
\begin{tabular}{@{}llr@{}}
\toprule Item & Détail & Pts\\\midrule
Q3.1–Q3.4 & déterminisme, yield, réduplication & 2\\
E3.1 & \texttt{run}/\texttt{accepts} du BUTA & 1\\
E3.2 & automate morphologique + \texttt{classify} & 0,5\\
E3.3 & AttackDecomposer & 0,5\\
E3.4 & produit (intersection) & 0,5\\
E3.5 & hash-consing (intern + round-trip) & 0,5\\
\textbf{Bonus} & infixation ; minimisation BUTA ; P4.5 (produit $A\times A_{enc}$) & +3\\
\bottomrule
\end{tabular}
\end{center}
\begin{gate}Les apps \textbf{instancient} \texttt{formlang.tree.TreeAutomaton}/\texttt{Term}.
Toute réimplémentation perd les points d'intégration.\end{gate}
\clearpage
% ===================== JOUR 4 =====================
\section{Jour 4 — Calculabilité : machine de Turing \& machine universelle}
\emph{Modules : \texttt{formlang/turing.py}, \texttt{utm.py} — App : \texttt{apps/mtu} —
Tests : \texttt{test\_turing.py}, \texttt{test\_calc.py} — Barème : 4 pts (+1).}

\subsection{Contexte}
On lève la dernière limite : une mémoire \textbf{libre}, lisible \emph{et} ré-inscriptible
n'importe où — la machine de Turing. Puis l'idée vertigineuse : \textbf{une seule} machine
peut simuler \textbf{toutes} les autres si on lui donne leur description. C'est la machine
\textbf{universelle}, ancêtre du processeur et de l'interpréteur.

\subsection{Notions nouvelles}
\begin{objetformel}{la machine de Turing}
\champ{Intuition.} Le modèle \textbf{ultime} : un ruban infini que l'on peut relire et
réécrire librement. Il capture « tout ce qui est mécaniquement calculable ».

\champ{Définition formelle.} $M=(Q,\Sigma,\Gamma,\delta,q_0,\text{\textvisiblespace},F)$ avec
$\delta:Q\times\Gamma\to Q\times\Gamma\times\{L,R,S\}$. Elle reconnaît les langages
\textbf{récursivement énumérables} et calcule les fonctions \textbf{calculables}.

\champ{Classe de problèmes.} \textbf{Tous} les problèmes de calcul : décision (avec la
frontière \emph{décidable}/\emph{indécidable}) et calcul de fonctions. C'est la limite du
calculable (thèse de Church--Turing).

\champ{Dans le projet.} Réaliser l'arithmétique unaire (\texttt{ADD}, \texttt{SUB}) comme de
\textbf{vraies} tables de transition, exécutées pas à pas.
\end{objetformel}

\begin{objetformel}{la machine universelle}
\champ{Intuition.} Une machine de Turing qui \textbf{lit la description d'une autre} machine
et l'exécute : un \textbf{interpréteur}.

\champ{Définition formelle.} $U(\langle M\rangle\,\#\#\,w)$ se comporte comme $M(w)$
(principe du \emph{programme enregistré}).

\champ{Classe de problèmes.} La \textbf{simulation universelle} — fondement de l'ordinateur
programmable — et la \textbf{frontière de l'indécidabilité} (le problème de l'arrêt).

\champ{Dans le projet.} Exécuter la calculatrice \textbf{par} $U$
(\texttt{addition\_via\_utm}) : programme-donnée, pas de boucle de ruban réécrite.
\end{objetformel}

\subsubsection{La machine de Turing}
$M=(Q,\Sigma,\Gamma,\delta,q_0,\text{\textvisiblespace},F)$ avec
$\delta:Q\times\Gamma\to Q\times\Gamma\times\{L,R,S\}$. À chaque pas : « dans l'état $q$, je
lis $a$ $\to$ j'écris $b$, je bouge de $d$, je passe en $q'$ ». Différence avec la pile : ici
on peut \textbf{relire et réécrire} le passé.
\begin{center}
\begin{tikzpicture}[start chain=going right, node distance=0mm]
  \foreach \c/\i in {\_/-1,1/0,1/1,+/2,1/3,\_/4,\_/5} {\node[tapecell, on chain] (c\i) {\c};}
  \node[above=5mm of c0, iutorange!85!black] (tete) {\small\bfseries tête, état \texttt{q0}};
  \draw[tapehead] (tete) -- (c0);
\end{tikzpicture}
\end{center}

\subsubsection{Une opération \emph{est} une machine : trace de \texttt{ADD} sur \texttt{11+1}}
\begin{exemple}[\texttt{2+1 = 3}, pas à pas]
\begin{center}\small\ttfamily
\begin{tabular}{@{}clll@{}}
\toprule \rmfamily pas & \rmfamily état & \rmfamily ruban & \rmfamily action\\\midrule
0 & q0 & [1]1+1 & lit 1 $\to$ R\\
1 & q0 & 1[1]+1 & lit 1 $\to$ R\\
2 & q0 & 11[+]1 & \rmfamily le \texttt+ devient \texttt1, R, passe q1\\
3 & q1 & 111[1] & lit 1 $\to$ R\\
4 & q1 & 1111[\_] & \rmfamily blanc : demi-tour, q2\\
5 & q2 & 111[1]\_ & \rmfamily efface le dernier 1, qf\\
6 & qf & 111[\_] & \rmfamily \bfseries accepté ; 111 = 3\\
\bottomrule
\end{tabular}
\end{center}
\texttt+$\to$\texttt1 \emph{fusionne} les deux blocs, puis on retire \textbf{un} \texttt1.
\end{exemple}

\subsubsection{Lire une table par \emph{phases} : la soustraction}
\texttt{SUB} apparie un \texttt1 de $m$ (marqué \texttt X) avec un \texttt1 de $n$ (barré,
droite$\to$gauche), jusqu'à épuiser $m$, puis nettoie.
\begin{center}
\begin{tikzpicture}[node distance=20mm and 22mm, on grid]
  \node[etat] (q0) {q0}; \node[etat, right=of q0] (q1) {q1};
  \node[etat, right=of q1] (q2) {q2}; \node[etat, below=of q1] (q3) {q3};
  \node[etat, right=22mm of q2] (q4) {q4}; \node[etatfin, right=of q4, accepting] (qf) {qf};
  \path[fleche]
    (q0) edge node[above]{\scriptsize \texttt{-} : R} (q1)
    (q1) edge node[above]{\scriptsize \texttt{1$\to$X}, L} (q2)
    (q2) edge[bend left=15] node[below,pos=.4]{\scriptsize \texttt{-} : L} (q3)
    (q3) edge[bend left=15] node[above,pos=.6]{\scriptsize \texttt{1$\to$X}, R} (q0)
    (q1) edge[bend left=20] node[above]{\scriptsize \texttt{\_} : $m$ épuisé} (q4)
    (q4) edge node[above]{\scriptsize nettoie} (qf);
\end{tikzpicture}
\end{center}
\textbf{Trois points délicats} (à expliquer au rapport) : \texttt{(q0,X)$\to$R} et
\texttt{(q1,X)$\to$R} (gérer ses propres marques) ; barrage \textbf{droite$\to$gauche}
(sortie contiguë) ; fin de $m$ par un \textbf{blanc} (le ruban compte, pas un compteur).

\subsubsection{Codage $\langle M\rangle$ et machine universelle}
Une MT n'est qu'une \textbf{table finie} : on l'écrit comme une \textbf{chaîne}
$\langle M\rangle$ (\textbf{injective}, \textbf{décodable}). $U$ prend
$\langle M\rangle\,\#\#\,w$ et \textbf{simule} $M$ (programme enregistré). Un cycle :
\begin{center}
\begin{tikzpicture}[noeudc/.style={draw=iutblue, very thick, rounded corners=3pt,
    fill=iutblue!6, minimum height=8mm, inner xsep=5pt, font=\small\bfseries, text=iutblue},
    fl/.style={-{Stealth[length=2mm]}, very thick, iutorange!85!black}]
  \node[noeudc] (lire) {1. Lire le symbole simulé};
  \node[noeudc, right=12mm of lire] (chercher) {2. Chercher dans $\langle M\rangle$};
  \node[noeudc, below=8mm of chercher] (ecrire) {3. Écrire};
  \node[noeudc, left=12mm of ecrire] (bouger) {4. Déplacer la tête};
  \node[noeudc, below=8mm of lire] (maj) {5. Mettre à jour l'état};
  \draw[fl] (lire)--(chercher); \draw[fl] (chercher)--(ecrire);
  \draw[fl] (ecrire)--(bouger); \draw[fl] (bouger)--(maj); \draw[fl] (maj)--(lire);
\end{tikzpicture}
\end{center}
\begin{notion}[Church--Turing, composition, limite]
Thèse de Church--Turing : pas un théorème, étayée par l'équivalence de modèles
($\lambda$-calcul, fonctions récursives, RAM). On obtient $\times,\div$ par
\textbf{composition} de MT (ce que fait $U$). \emph{Revers} : \textsc{halt} est
\textbf{indécidable} (Turing 1936) ; par réduction, l'arrêt de $U$ l'est aussi.
\emph{Universel $\neq$ omniscient.}
\end{notion}

\subsection{A. Questions théoriques \normalfont(à rédiger) — 2 pts}
\begin{description}[leftmargin=1.8em,style=nextline]
  \item[Q4.1 (0,5) — encodage.] Détailler $\langle M\rangle$ (numéroter états/symboles,
    linéariser $\delta$ en quintuplets). \emph{Livrable} : le schéma + preuve d'injectivité
    et de décodabilité.
  \item[Q4.2 (0,5) — la machine universelle.] Décrire $U$, l'invariant maintenu sur le
    ruban, et un cycle complet. \emph{Livrable} : l'invariant + l'argument de correction
    (récurrence sur le nombre de pas).
  \item[Q4.3 (0,5) — surcoût (\textbf{sourcé}).] Multi-ruban $O(t\cdot|\langle M\rangle|)$ ;
    multi$\to$mono $\le$ quadratique ; citer Hennie--Stearns (ralentissement logarithmique).
    \emph{Livrable} : les bornes \textbf{avec référence} (aucune constante « de mémoire »).
  \item[Q4.4 (0,5) — indécidabilité.] Réduction depuis \textsc{halt} montrant que « $U$
    s'arrête sur $\langle M\rangle\#\#w$ » est indécidable ; commenter « universel $\neq$
    omniscient ».
\end{description}

\subsection{B. Exercices de programmation — 2 pts}
\begin{description}[leftmargin=1.8em,style=nextline]
  \item[E4.1 (0,5) — MT générique.] \texttt{TuringMachine.run} (ruban \texttt{dict}, arrêt
    sur état final / transition absente, option \texttt{trace}). \emph{Test} :
    \texttt{test\_trace\_disponible}.
  \item[E4.2 (0,5) — machine universelle.] \texttt{encode}/\texttt{decode} (réciproques
    exactes) + \texttt{UniversalTM.run}. \emph{Tests} : \texttt{test\_round\_trip\_encodage},
    \texttt{test\_utm\_simule\_comme\_la\_machine}.
  \item[E4.3 (0,75) — calculatrice.] Tables \texttt{ADD}/\texttt{SUB} + \texttt{Calculatrice}
    (\texttt+,\texttt- par les MT ; \texttt*,\texttt/ par composition ; \texttt{chainer}).
    \emph{Tests} : \texttt{test\_add\_sub\_machines\_directes}, \texttt{test\_calculatrice\_exhaustif},
    \texttt{test\_chainage}, \texttt{test\_division\_par\_zero}.
  \item[E4.4 (0,25) — interpréteur universel.] \texttt{UniversalInterpreter.run(M,w)} qui
    encode $M$ puis \textbf{délègue} à \texttt{UniversalTM} ; \texttt{addition\_via\_utm}.
    \emph{Test} : \texttt{test\_interpreteur\_universel}.
\end{description}
\begin{notion}[Pièges fréquents]
\textbullet\ \texttt{run} : ne pas garantir l'arrêt (prévoir \texttt{max\_steps} et lever
en cas de boucle).\\
\textbullet\ \texttt{SUB} : oublier \texttt{(q0,X)$\to$R}/\texttt{(q1,X)$\to$R} $\Rightarrow$
blocage au 2\textsuperscript{e} tour.\\
\textbullet\ \texttt{encode}/\texttt{decode} : sortie non triée $\Rightarrow$ round-trip
non déterministe.
\end{notion}
\begin{center}\small
\begin{tabular}{@{}llr@{}}
\toprule Item & Détail & Pts\\\midrule
Q4.1–Q4.4 & encodage, cycle $U$, surcoût \textbf{sourcé}, indécidabilité & 2\\
E4.1 & MT générique & 0,5\\
E4.2 & encode/decode + \texttt{UniversalTM} & 0,5\\
E4.3 & calculatrice (tests exhaustifs) & 0,75\\
E4.4 & interpréteur universel & 0,25\\
\textbf{Bonus} & table \textbf{monolithique} de $\times$ validée par le simulateur & +1\\
\bottomrule
\end{tabular}
\end{center}
\begin{gate}\texttt{apps/mtu} ne réécrit aucune boucle de MT : il décrit des \textbf{tables}
et les fait tourner via \texttt{formlang.turing}/\texttt{utm}. Surcoût \textbf{référencé}.\end{gate}
\clearpage

% ===================== JOUR 5 =====================
\section{Jour 5 — Intégration \& Myhill--Nerode}
\emph{Modules : \texttt{pipeline.py}, \texttt{formlang/myhill\_nerode.py} — Tests :
\texttt{test\_pipeline.py}, \texttt{test\_cfg\_nerode.py} — Barème : 3 pts (+1).}

\subsection{Contexte}
Deux visages : \textbf{assembler} (les étages coopèrent dans un \texttt{pipeline}) et
\textbf{boucler la boucle théorique} — revenir sur « combien d'états faut-il vraiment ? »
avec la \textbf{congruence de Myhill--Nerode}.

\subsection{Notions nouvelles}
\begin{objetformel}{la congruence de Myhill--Nerode}
\champ{Intuition.} Regrouper les mots qui « ont le même \textbf{avenir} » vis-à-vis du
langage : une fois lus, ils mènent aux mêmes acceptations futures.

\champ{Définition formelle.} $u\approx_{L}v$ ssi, pour tout suffixe $w$,
$uw\in L\Leftrightarrow vw\in L$. C'est une équivalence \textbf{invariante à droite}, et
$L$ est régulier \textbf{ssi} son \textbf{indice} (nombre de classes) est fini.

\champ{Classe de problèmes.} La \textbf{caractérisation} des langages réguliers et la
\textbf{minimisation} : l'indice égale le nombre d'états de l'AFD minimal ; identifier les
états équivalents, c'est identifier les classes.

\champ{Dans le projet.} Retrouver les \textbf{3 classes} de $L_1$ (= 3 états de l'AFD minimal
du jour~1) et refermer ainsi le fil rouge théorique.
\end{objetformel}

\subsubsection{Intégration : deux chaînes, un même cœur}
\begin{center}
\begin{tikzpicture}[etape/.style={draw=iutblue,very thick,rounded corners=3pt,fill=iutblue!6,
    minimum height=7mm,inner xsep=5pt,font=\footnotesize\bfseries,text=iutblue,align=center},
    fl/.style={-{Stealth[length=2mm]},very thick,iutorange!85!black}, node distance=8mm]
  \node[etape] (fst) {FST\\normalise}; \node[etape,right=of fst] (dfa) {AFD\\facteur \texttt{or}};
  \node[etape,right=of dfa] (pda) {PDA\\délimiteurs};
  \node[left=6mm of fst,font=\scriptsize\itshape,text=iutgrey] {\texttt{analyze\_word}};
  \draw[fl] (fst)--(dfa); \draw[fl] (dfa)--(pda);
  \node[etape,below=9mm of fst] (disc) {découverte\\affixes};
  \node[etape,right=of disc] (seg) {segmentation}; \node[etape,right=of seg] (buta) {BUTA\\classify};
  \node[left=6mm of disc,font=\scriptsize\itshape,text=iutgrey] {\texttt{analyze\_morpho}};
  \draw[fl] (disc)--(seg); \draw[fl] (seg)--(buta);
\end{tikzpicture}
\end{center}

\subsubsection{La congruence de Myhill--Nerode}
$u\approx_{L}v$ si, \textbf{pour tout} suffixe $w$, $uw\in L\Leftrightarrow vw\in L$ :
l'\emph{avenir} après $u$ ou $v$ est identique. Équivalence \textbf{invariante à droite}.
\begin{notion}[Le théorème — le fil rouge se referme]
$L$ est régulier \textbf{ssi} $\approx_{L}$ a un nombre \textbf{fini} de classes, et ce
nombre (l'\textbf{indice}) est \textbf{exactement} le nombre d'états de l'AFD \textbf{minimal}.
Minimiser (jour~1), c'était calculer ces classes.
\end{notion}

\subsubsection{En pratique : signatures sur suffixes témoins}
Avec $S=\{\varepsilon, r, or\}$, la \textbf{signature} d'un mot est
$\big(\mathbf 1[w s\in L_1]\big)_{s\in S}$ :
\begin{exemple}[Trois signatures = trois classes]
\begin{center}\small
\begin{tabular}{@{}lccc>{\bfseries}c@{}}
\toprule mot $w$ & $w\varepsilon$ & $wr$ & $wor$ & classe\\\midrule
\texttt{""} & 0 & 0 & 1 & A\\ \texttt{a} & 0 & 0 & 1 & A\\
\texttt{o} & 0 & 1 & 1 & B\\ \texttt{ao} & 0 & 1 & 1 & B\\
\texttt{or} & 1 & 1 & 1 & C\\ \texttt{aor} & 1 & 1 & 1 & C\\
\bottomrule
\end{tabular}
\end{center}
Signatures $(0,0,1),(0,1,1),(1,1,1)$ : exactement \textbf{3 classes} $\Rightarrow$ 3 états.
\end{exemple}
\begin{center}
\begin{tikzpicture}[classe/.style={draw=iutblue,very thick,rounded corners=4pt,fill=iutblue!5,
    minimum width=33mm,minimum height=14mm,align=center,font=\small}]
  \node[classe] (A) {\textbf{$C_A$}\\ \texttt{""},\ \texttt{a}\\[-1pt]{\footnotesize\bfseries\color{iutorange!85!black}état A}};
  \node[classe,right=8mm of A] (B) {\textbf{$C_B$}\\ \texttt{o},\ \texttt{ao}\\[-1pt]{\footnotesize\bfseries\color{iutorange!85!black}état B}};
  \node[classe,right=8mm of B] (C) {\textbf{$C_C$}\\ \texttt{or},\ \texttt{aor}\\[-1pt]{\footnotesize\bfseries\color{iutorange!85!black}état C}};
  \draw[-{Stealth[length=2mm]},thick,iutblue!70] (A)--node[above,font=\scriptsize]{$o$}(B);
  \draw[-{Stealth[length=2mm]},thick,iutblue!70] (B)--node[above,font=\scriptsize]{$r$}(C);
\end{tikzpicture}
\end{center}

\subsection{A. Questions théoriques \normalfont(à rédiger) — 1 pt}
\begin{description}[leftmargin=1.8em,style=nextline]
  \item[Q5.1 (0,5)] Énoncer Myhill--Nerode et relier le nombre de classes de $L_1$ au
    nombre d'états de l'AFD minimal du jour~1. \emph{Livrable} : énoncé + le lien chiffré (3).
  \item[Q5.2 (0,5)] Expliquer le critère « indistinguables par tous les suffixes témoins »
    et son lien avec la \textbf{fusion d'états} à la minimisation.
\end{description}

\subsection{B. Exercices de programmation — 2 pts}
\begin{description}[leftmargin=1.8em,style=nextline]
  \item[E5.1 (1) — pipeline.] \texttt{analyze\_word} (FST$\to$AFD$\to$PDA) et
    \texttt{analyze\_morpho} (\texttt{discover}$\to$\texttt{segment\_to\_tree}$\to$\texttt{classify}).
    \emph{Tests} : \texttt{test\_pipeline\_word}, \texttt{test\_pipeline\_morpho}.
  \item[E5.3 (1) — classes de Nerode.] \texttt{nerode\_classes(accepts, words, suffixes)} :
    regrouper par signature \texttt{tuple(accepts(w+s) for s in suffixes)}. \emph{Test} :
    \texttt{test\_nerode\_trois\_classes}.
\end{description}
\begin{exemple}[Sortie de référence]\ttfamily
\$ pytest -q\\ 28 passed\\[2pt]
\$ python pipeline.py --word 4or\\
\hspace*{2em}normalisé(FST) : aor\quad facteur\_or(AFD) : True\quad délimiteurs\_ok(PDA) : True
\end{exemple}
\begin{center}\small
\begin{tabular}{@{}llr@{}}
\toprule Item & Détail & Pts\\\midrule
Q5.1–Q5.2 & Myhill--Nerode + lien minimisation & 1\\
E5.1 & pipeline (mot + morpho) & 1\\
E5.3 & classes de Nerode & 1\\
\textbf{Bonus} & mesure hash-consing sur 10\,000 mots & +1\\
\bottomrule
\end{tabular}
\end{center}
\begin{gate}Le pipeline \textbf{appelle les apps} (qui instancient le cœur) ; il ne
reprogramme aucun automate.\end{gate}
\clearpage

% ===================== BIBLIOGRAPHIE =====================
\section{Bibliographie (vérifiée)}
Les références marquées \textbf{(libre)} sont disponibles gratuitement.

\paragraph{Automates finis \& langages.} M.~Sipser, \emph{Introduction to the Theory of
Computation}, 3\textsuperscript{e}~éd., Cengage, 2012. \textbullet\ J.~Hopcroft, R.~Motwani,
J.~Ullman, \emph{Introduction to Automata Theory, Languages, and Computation},
3\textsuperscript{e}~éd., Pearson, 2006.

\paragraph{Automates d'arbres.} H.~Comon \emph{et al.}, \emph{Tree Automata Techniques and
Applications (TATA)}, 2007 \textbf{(libre : \texttt{tata.gforge.inria.fr})}. \textbullet\
F.~Gécseg, M.~Steinby, \emph{Tree Automata}, 1984 ; rééd.\ \textbf{(libre : arXiv:1509.06233,
2015)}. \textbullet\ J.~W.~Thatcher, J.~B.~Wright, \emph{Math.\ Systems Theory} 2(1):57–81,
1968 \emph{(fondateur)}. \textbullet\ J.~Doner, \emph{JCSS} 4(5):406–451, 1970.

\paragraph{Applications des arbres.} M.~Murata \emph{et al.}, \emph{ACM TOIT} 5(4), 2005
(schémas XML). \textbullet\ A.~V.~Aho, M.~Ganapathi, S.~Tjiang, \emph{ACM TOPLAS} 11(4),
1989 (filtrage d'arbres). \textbullet\ A.~Bouajjani \emph{et al.}, \emph{ENTCS} 149(1), 2006.

\paragraph{Morphologie computationnelle.} K.~Koskenniemi, \emph{Two-Level Morphology},
Helsinki, 1983. \textbullet\ K.~R.~Beesley, L.~Karttunen, \emph{Finite State Morphology},
CSLI, 2003. \textbullet\ Z.~S.~Harris, \emph{Language} 31(2):190–222, 1955. \textbullet\
J.~Goldsmith, \emph{Computational Linguistics} 27(2):153–198, 2001. \textbullet\ E.~Roche,
Y.~Schabes (dir.), \emph{Finite-State Language Processing}, MIT Press, 1997.

\paragraph{Réduplication \& limites.} C.~Culy, « The complexity of the vocabulary of
Bambara », \emph{Linguistics and Philosophy} 8:345–351, 1985. \textbullet\ B.~Roark,
R.~Sproat, \emph{Computational Approaches to Morphology and Syntax}, Oxford, 2007.

\paragraph{Hash-consing.} J.-C.~Filliâtre, S.~Conchon, « Type-safe modular hash-consing »,
\emph{ML Workshop}, 2006. \textbullet\ J.~Daciuk \emph{et al.}, \emph{Computational
Linguistics} 26(1), 2000.

\paragraph{Calculabilité.} A.~M.~Turing, « On Computable Numbers\dots », \emph{Proc.\
London Math.\ Soc.}, s2-42:230–265, 1936-37. \textbullet\ F.~C.~Hennie, R.~E.~Stearns,
« Two-Tape Simulation of Multitape Turing Machines », \emph{J.\ ACM} 13(4):533–546, 1966.
\textbullet\ S.~Arora, B.~Barak, \emph{Computational Complexity: A Modern Approach},
Cambridge, 2009.

\end{document}
